\scanpage{289}
\mspage{30}

\[
\begin{aligned}
\mathrm{Hom}_{\underline{\mathrm{Typ}}}(\tau,\tau')&=\mathrm{Hom}_{\mathrm{Cat}}(I_\tau,I_{\tau'})\\
\mathrm{Hom}_{\mathrm{Ty}}(\tau,\tau')&=\mathrm{Hom}^{\mathrm{int}}_{\mathrm{Cat}}(I_\tau,I_{\tau'})
\end{aligned}
\]
(hom. respectant la structure d'intervalle)

On aura
\[
\underline{\mathrm{C}\deleted{\unsure{\mathrm{H}}}\mathrm{h}}(X)\simeq \underline{\mathrm{CH}}(X)\mathbin{\times}_{\underline{\mathrm{Typ}}}\underline{\mathrm{Ty}},
\]
i.e. $\underline{\mathrm{Ch}}(X)$ est la \uline{restriction} de la catégorie fibrée $\underline{\mathrm{CH}}(X)$ sur $\underline{\mathrm{Typ}}$, à la sous-catégorie (non pleine) $\underline{\mathrm{Ty}}$ de $\underline{\mathrm{Typ}}$.

\uline{NB} Les catégories fibrées $\underline{\mathrm{CH}}(X),\underline{\mathrm{Ch}}(X)$ sur $\underline{\mathrm{Typ}},\underline{\mathrm{Ty}}$ dépendent fonctoriellement de $X$, variant dans $(\mathrm{Cat})$, comme il se doit. Comme de juste, on s'y intéresse surtout quand $X$ \uline{n'est pas} dans $(\mathrm{Cat})$, i.e. quand $X$ est une ``grosse'' catégorie (telle que $(\mathrm{Cat})$ elle-même). Je me dispense \deleted{\unsure{de}}d'expliciter ici les questions pertinentes d'Univers.

Deux tests essentiels, pour s'assurer si la notion est bonne

a) La théorie des $\unsure{M|\mathcal M}\Sigma^{-1}$

b) La composition des chemins) doit bien marcher.
% archive sheet number: 288
